\section{Introduction}\label{sec:1}\label{introduction}

The difficult model comparison need not be the pair with the smallest performance difference. If differences up to two percentage points are practically equivalent, a difference near zero lies inside a decision region; a difference near two points lies on its edge. Early-stopping cost should therefore be studied against distance to that edge, rather than summarized only over arbitrary model pairs.

We ask: \textbf{when can certified model comparison stop substantially early, and when is the remaining uncertainty about unseen items fundamentally consequential?} We condition on a fixed finite benchmark and fixed paired scores, so randomness comes from the evaluation design. A correct decision means practical superiority beyond a predeclared margin or affirmative equivalence within it.

Sequential equivalence and without-replacement confidence sequences supply the inferential tools \citep{r1,r2}. Sequential model evaluation already studies stopping cost \citep{r3}, and election and financial audits certify assertions about fixed finite populations \citep{r18,r19,r20}. Building on this literature, we separate the query cost imposed by label-changing alternatives from the cost of a chosen inference procedure.

The main theoretical result extends the finite-vector inclusion argument to nonconstant near-boundary populations through their available coordinate changes. It recovers a near-census lower bound at one-item resolution. A standard concentration argument supplies an explicit upper bound away from the boundary. The necessary and sufficient bounds operate at different scales. The main empirical result is a boundary-dependent cost landscape, with substantial differences between low- and high-variance populations at the same gap. Figure~\ref{fig:boundary} places controlled populations beside both random public pairs and a deliberately difficult cohort. The theory establishes an obstruction when few unseen changes can alter the label; the replays measure how much of the benchmark specific certified procedures actually inspect. Their agreement is qualitative outside the resolution covered by the lower bound.
